% COMP0197 C2 — Progress Presentation & Design Defence (8 min + Q&A)
% 4 speakers (one per required discussion point) + 3 support + 2 Q&A defenders
\documentclass[11pt,a4paper]{article}

\usepackage[T1]{fontenc}
\usepackage[margin=2.2cm]{geometry}
\usepackage{amsmath,amssymb}
\usepackage{booktabs}
\usepackage{hyperref}
\usepackage{xcolor}
\usepackage{enumitem}
\usepackage{tikz}
\usetikzlibrary{shapes.geometric, arrows.meta, positioning}
\usepackage{tcolorbox}
\tcbuselibrary{skins,breakable}

\newtcolorbox{speakerbox}[2][]{
  colback=#2!5, colframe=#2!70!black, fonttitle=\bfseries\large,
  title={#1}, breakable, left=4pt, right=4pt, top=4pt, bottom=4pt
}
\newtcolorbox{talkingpoints}[1][]{
  colback=gray!5, colframe=gray!50, fonttitle=\bfseries,
  title={Talking Points}, left=3pt, right=3pt, top=2pt, bottom=2pt, #1
}
\newtcolorbox{anticipatedqa}[1][]{
  colback=red!3, colframe=red!40, fonttitle=\bfseries,
  title={Anticipated Q\&A (for Defenders)}, left=3pt, right=3pt, top=2pt, bottom=2pt, #1
}

\title{\textbf{ProbTSLANet: C2 Presentation Script}\\[4pt]
\large COMP0197 --- Progress Presentation \& Design Defence ($\sim$8 min + Q\&A)\\[3pt]
\normalsize 7 Presenting Speakers + 2 Q\&A Defenders}
\author{Group Members}
\date{}

\begin{document}
\maketitle
\thispagestyle{empty}

% ============================================================================
\section*{C2 Requirements (from spec)}
% ============================================================================

The spec defines \textbf{four} key discussion points that must be covered:
\begin{enumerate}[leftmargin=*]
    \item \textbf{Design and Rationale}: Justification for the selected problem, dataset choice and the specific high-level architecture used.
    \item \textbf{Implementation Plan}: A roadmap of how the system components (data loader, model, training loop and evaluation) are being integrated.
    \item \textbf{Current Progress}: Demonstration of functional code segments, such as a working data pipeline or initial training benchmarks.
    \item \textbf{Problem Solving}: How the group has addressed technical hurdles like data imbalance, slow convergence or hardware limitations.
\end{enumerate}

\noindent The group must \emph{demonstrate and defend through Q\&A a clear path to a working system}. No code is submitted at this stage.

% ============================================================================
\section*{Role Assignment ($\sim$8 min presentation + open Q\&A)}
% ============================================================================

\begin{center}
\renewcommand{\arraystretch}{1.25}
\begin{tabular}{@{}cllc@{}}
\toprule
\textbf{\#} & \textbf{Role} & \textbf{C2 Discussion Point} & \textbf{Time} \\
\midrule
1 & \textcolor{blue!70!black}{\textbf{Speaker 1}} & Design \& Rationale (problem + dataset) & $\sim$1.5 min \\
2 & \textcolor{orange!80!black}{\textbf{Speaker 2}} & Design \& Rationale (architecture) & $\sim$1.5 min \\
3 & \textcolor{teal!80!black}{\textbf{Speaker 3}} & Implementation Plan (components roadmap) & $\sim$1.5 min \\
4 & \textcolor{violet!80!black}{\textbf{Speaker 4}} & Current Progress (working demo + benchmarks) & $\sim$1.5 min \\
5 & \textcolor{red!70!black}{\textbf{Speaker 5}} & Problem Solving (technical hurdles) & $\sim$1 min \\
6 & \textcolor{green!50!black}{\textbf{Speaker 6}} & Problem Solving (what's next) & $\sim$0.5 min \\
7 & \textcolor{brown!80!black}{\textbf{Speaker 7}} & Wrap-up + transition to Q\&A & $\sim$0.5 min \\
\midrule
8 & \textcolor{gray!70!black}{\textbf{Defender A}} & Q\&A --- Theory \& design justification & as needed \\
9 & \textcolor{gray!70!black}{\textbf{Defender B}} & Q\&A --- Code \& implementation details & as needed \\
\bottomrule
\end{tabular}
\end{center}

\newpage

% ============================================================================
% SPEAKER 1 — Design & Rationale: Problem + Dataset
% ============================================================================
\begin{speakerbox}[Speaker 1 --- Design \& Rationale: Problem + Dataset \textnormal{($\sim$1.5 min)}]{blue}

\textbf{C2 requirement addressed:} \textit{Justification for the selected problem and dataset choice.}

\begin{talkingpoints}
\begin{itemize}[leftmargin=*]
    \item \textbf{Problem}: Standard forecasting models output a single number --- but real decisions need \emph{confidence}. We build a system that outputs a \textbf{probability distribution} $\mathcal{N}(\mu, \sigma^2)$ for every future timestep.
    \item \textbf{Two types of uncertainty} (one sentence each):
    \begin{itemize}
        \item \emph{Aleatoric}: irreducible data noise (measurement error, chaos).
        \item \emph{Epistemic}: model uncertainty, reducible with more data.
    \end{itemize}
    Our system decomposes and reports both.
    \item \textbf{Dataset}: Weather benchmark --- 52K hourly observations, 21 meteorological variables, from a German weather station. Standard in the forecasting literature (Autoformer, PatchTST).
    \item \textbf{Why weather?} Strong periodic patterns (ideal for spectral modelling) + genuine noise (ideal for testing uncertainty quantification). Chronological 70/10/20 split, no leakage.
\end{itemize}
\end{talkingpoints}

\begin{anticipatedqa}
\textbf{Q: Why only one dataset?} --- We prioritise depth of analysis (7-config ablation $\times$ 2 horizons = 14 experiments) over breadth. Weather is a standard benchmark allowing comparison with published results. \\
\textbf{Q: Why not finance/medical data?} --- Weather has well-understood physics making uncertainty results interpretable. The framework generalises to other domains.
\end{anticipatedqa}

\end{speakerbox}

\vspace{6pt}

% ============================================================================
% SPEAKER 2 — Design & Rationale: Architecture
% ============================================================================
\begin{speakerbox}[Speaker 2 --- Design \& Rationale: Architecture \textnormal{($\sim$1.5 min)}]{orange}

\textbf{C2 requirement addressed:} \textit{Justification for the specific high-level architecture.}

\begin{talkingpoints}
\begin{itemize}[leftmargin=*]
    \item \textbf{Backbone}: TSLANet (ICML 2024) --- lightweight, patch-based, combines spectral filtering with convolutional gating. We extend it into \textbf{ProbTSLANet}.
    \item \textbf{High-level flow} (show diagram):
    \begin{center}
    \small Input $\to$ Instance Norm $\to$ Patch \& Embed $\to$ 2$\times$ TSLANet Blocks $\to$ Flatten $\to$ $\begin{cases} \mu\text{-head} \\ \log\sigma^2\text{-head} \end{cases}$ $\to$ $\mathcal{N}(\mu,\sigma^2)$
    \end{center}
    \item \textbf{Key components} (one line each):
    \begin{itemize}
        \item \emph{Adaptive Spectral Block}: FFT $\to$ learnable frequency weights + energy-based mask $\to$ iFFT. Separates trend from noise in frequency domain.
        \item \emph{Interacting Conv Block}: Two parallel convolutions with multiplicative gating --- captures local temporal patterns.
        \item \emph{Dual output heads}: $\mu$-head for the mean, $\log\sigma^2$-head for learned variance (heteroscedastic Gaussian).
    \end{itemize}
    \item \textbf{Why TSLANet?} Spectral block matches weather periodicity; channel-independent design scales to 21 variables; lightweight ($\sim$50K params, trains in minutes).
\end{itemize}
\end{talkingpoints}

\begin{anticipatedqa}
\textbf{Q: Why not Transformer / LSTM?} --- TSLANet's spectral block directly learns frequency representations ideal for periodic weather data, avoids quadratic attention cost, and the original paper shows competitive results on Weather. \\
\textbf{Q: Why depth 2?} --- Diminishing returns beyond 2 for this dataset size; deeper models overfit on 36K samples.
\end{anticipatedqa}

\end{speakerbox}

\newpage

% ============================================================================
% SPEAKER 3 — Implementation Plan
% ============================================================================
\begin{speakerbox}[Speaker 3 --- Implementation Plan \textnormal{($\sim$1.5 min)}]{teal}

\textbf{C2 requirement addressed:} \textit{A roadmap of how the system components (data loader, model, training loop and evaluation) are being integrated.}

\begin{talkingpoints}
\begin{itemize}[leftmargin=*]
    \item \textbf{Codebase structure}: Two self-contained scripts, no external model libraries:
    \begin{itemize}
        \item \texttt{train.py} --- data loading, model definition, pretraining, fine-tuning, saves weights + config.
        \item \texttt{test.py} --- loads model, runs inference, computes metrics, generates plots.
    \end{itemize}
    \item \textbf{Component integration roadmap}:
    \begin{enumerate}[leftmargin=*]
        \item \textbf{Data loader}: \texttt{Dataset\_Custom} class --- reads CSV, z-score normalisation (fitted on train only), sliding-window sampling, time feature extraction. $\checkmark$ Done.
        \item \textbf{Model}: \texttt{ProbabilisticTSLANet} --- patching, 2$\times$ TSLANet blocks (ASB+ICB), dual output heads. $\checkmark$ Done.
        \item \textbf{Training loop}: Phase 1 (self-supervised masked patch pretraining, 5 epochs) $\to$ Phase 2 (supervised fine-tuning with NLL or MSE loss, 15 epochs). AdamW, LR scheduling, gradient clipping. $\checkmark$ Done.
        \item \textbf{Evaluation}: Point metrics (MSE, MAE) + probabilistic metrics (NLL, CRPS, calibration, sharpness) + uncertainty decomposition + visualisation. $\checkmark$ Done.
    \end{enumerate}
    \item \textbf{Ablation automation}: \texttt{run\_all.sh} runs all 14 experiments (7 configs $\times$ 2 horizons) end-to-end. Each experiment saves its own \texttt{config.json}, weights, metrics, and plots.
    \item \textbf{Framework}: Pure PyTorch (\texttt{comp0197-pt} environment). Only extra package: \texttt{matplotlib}.
\end{itemize}
\end{talkingpoints}

\begin{anticipatedqa}
\textbf{Q: Why single-file scripts?} --- Coursework requires standalone \texttt{train.py}/\texttt{test.py}. Embedding everything avoids import-path issues and ensures reproducibility. \\
\textbf{Q: How is reproducibility ensured?} --- Fixed seed (42), deterministic cuDNN, config saved per experiment, \texttt{instruction.pdf} with exact steps.
\end{anticipatedqa}

\end{speakerbox}

\vspace{6pt}

% ============================================================================
% SPEAKER 4 — Current Progress
% ============================================================================
\begin{speakerbox}[Speaker 4 --- Current Progress \textnormal{($\sim$1.5 min)}]{violet}

\textbf{C2 requirement addressed:} \textit{Demonstration of functional code segments, such as a working data pipeline or initial training benchmarks.}

\begin{talkingpoints}
\begin{itemize}[leftmargin=*]
    \item \textbf{All 14 experiments are complete}. Show summary of initial benchmarks:

    \begin{center}
    \footnotesize
    \setlength{\tabcolsep}{3pt}
    \begin{tabular}{@{}ll|cc|cc@{}}
    \toprule
    $H$ & Config & MSE$\downarrow$ & MAE$\downarrow$ & CRPS$\downarrow$ & Cal.\ Err.$\downarrow$ \\
    \midrule
    \multirow{3}{*}{96}
    & A1 (deterministic) & \textbf{.174} & \textbf{.215} & --- & --- \\
    & A3 (aleatoric) & .197 & .223 & \textbf{.185} & \textbf{.102} \\
    & A7 (full+auxMSE) & .183 & .217 & .194 & .139 \\
    \bottomrule
    \end{tabular}
    \end{center}

    \item \textbf{Key observation so far}: The Gaussian NLL head (A3) is the critical component --- produces well-calibrated intervals (CE = 0.102). MC Dropout alone (A2) fails completely (CE = 0.436, collapsed intervals).
    \item \textbf{Working demos to show if asked}: data pipeline loading weather CSV, training loss curves, prediction interval plots, calibration plots, uncertainty heatmaps.
    \item \textbf{Saved models}: 14 directories with weights, configs, metrics JSON, and PDF plots --- all reproducible.
\end{itemize}
\end{talkingpoints}

\begin{anticipatedqa}
\textbf{Q: Why is MSE worse for probabilistic models?} --- NLL loss optimises distributional accuracy, not point accuracy. A5/A7 add auxiliary MSE to recover most of the gap (0.183 vs.\ 0.174). \\
\textbf{Q: Show us the calibration plot?} --- (Have the saved PDF plots ready to display from \texttt{saved\_models/A3\_pl96/results/plots/}).
\end{anticipatedqa}

\end{speakerbox}

\newpage

% ============================================================================
% SPEAKER 5 — Problem Solving: Technical Hurdles
% ============================================================================
\begin{speakerbox}[Speaker 5 --- Problem Solving: Technical Hurdles \textnormal{($\sim$1 min)}]{red}

\textbf{C2 requirement addressed:} \textit{How the group has addressed technical hurdles.}

\begin{talkingpoints}
\begin{itemize}[leftmargin=*]
    \item \textbf{Hurdle 1 --- NLL loss instability}: Early training with NLL caused exploding gradients and the model ``cheating'' by inflating $\sigma^2$. \\
    \textbf{Solution}: (a) Conservative $\log\sigma^2$-head initialisation (weights $= 0$, bias $= -2$), (b) output clamping to $[-10, 10]$, (c) gradient clipping (max norm 4.0).

    \item \textbf{Hurdle 2 --- NLL hurting point accuracy}: Pure NLL sacrificed MSE to improve likelihood. \\
    \textbf{Solution}: Added auxiliary MSE loss ($\lambda = 0.3$) in A5/A7, recovering 7\% MSE while maintaining calibration. Also tried two-stage training (A6) but found it over-disperses.

    \item \textbf{Hurdle 3 --- MC Dropout alone is useless} (A2): Dropout variance from a deterministic model is near-zero, producing collapsed intervals. \\
    \textbf{Solution}: Recognised that aleatoric uncertainty \emph{must} be modelled explicitly via the Gaussian head; MC Dropout is only useful \emph{on top of} it.

    \item \textbf{Hurdle 4 --- Variance denormalisation}: Instance normalisation (RevIN) required careful variance scaling ($\log\sigma^2_{\text{out}} = \log\tilde{\sigma}^2 + 2\log s$) to avoid incorrect uncertainty in original scale.
\end{itemize}
\end{talkingpoints}

\end{speakerbox}

\vspace{6pt}

% ============================================================================
% SPEAKER 6 — Problem Solving: What's Next
% ============================================================================
\begin{speakerbox}[Speaker 6 --- What Remains \textnormal{($\sim$0.5 min)}]{green!50!black}

\begin{talkingpoints}
\begin{itemize}[leftmargin=*]
    \item \textbf{Remaining work for C3}:
    \begin{itemize}
        \item Finalise the comparison analysis across all 14 experiments.
        \item Polish visualisation (prediction intervals, calibration, uncertainty heatmaps).
        \item Write the LNCS report (Introduction, Methods, Experiments, Results, Discussion).
        \item Each member writes individual Part B (contribution, critique, GenAI audit).
    \end{itemize}
    \item \textbf{No major technical risks remain} --- all code runs, all experiments complete, all metrics computed.
\end{itemize}
\end{talkingpoints}

\end{speakerbox}

\vspace{6pt}

% ============================================================================
% SPEAKER 7 — Wrap-up
% ============================================================================
\begin{speakerbox}[Speaker 7 --- Wrap-up \textnormal{($\sim$0.5 min)}]{brown}

\begin{talkingpoints}
\begin{itemize}[leftmargin=*]
    \item One-line summary: \textbf{ProbTSLANet} extends TSLANet with Gaussian output heads and MC Dropout to produce calibrated probabilistic weather forecasts with decomposed aleatoric/epistemic uncertainty.
    \item The Gaussian NLL head is the key component; auxiliary MSE loss is the key practical fix.
    \item System is fully implemented and all 14 ablation experiments are complete.
    \item ``We're happy to take questions --- over to the panel.''
\end{itemize}
\end{talkingpoints}

\end{speakerbox}

\newpage

% ============================================================================
% Q&A DEFENDERS
% ============================================================================
\begin{speakerbox}[Defender A --- Q\&A: Theory \& Design Rationale]{gray}

\textbf{Likely questions and prepared answers:}
\begin{itemize}[leftmargin=*]
    \item \textbf{``Why Gaussian and not another distribution?''} --- Gaussian is the simplest continuous distribution with learnable mean and variance. Analytically tractable for NLL, CRPS. Richer options (mixture density, normalising flows) are future work.
    \item \textbf{``Is MC Dropout a valid Bayesian method?''} --- It approximates variational inference. Tends to \emph{underestimate} epistemic uncertainty vs.\ deep ensembles, but is computationally cheap (single model).
    \item \textbf{``Why is epistemic uncertainty so low?''} --- (1) Weather noise genuinely dominates, (2) MC Dropout with a small model underestimates it. Ensembles would give a stronger signal.
    \item \textbf{``What is the adaptive spectral mask actually doing?''} --- It computes per-frequency energy, thresholds against the median, and applies separate learnable weights to high-energy vs.\ low-energy frequencies. This lets the model treat periodic signals and noise differently.
    \item \textbf{``Is NLL a proper scoring rule?''} --- Yes, for the Gaussian family. CRPS is also strictly proper and more robust to outliers.
\end{itemize}
\end{speakerbox}

\vspace{6pt}

\begin{speakerbox}[Defender B --- Q\&A: Implementation \& Code Details]{gray}

\textbf{Likely questions and prepared answers:}
\begin{itemize}[leftmargin=*]
    \item \textbf{``How do you prevent data leakage?''} --- Chronological split (70/10/20), scaler fitted on train only, no shuffling at test time.
    \item \textbf{``How is MC Dropout implemented?''} --- \texttt{model.eval()} then \texttt{m.train()} only for \texttt{nn.Dropout} layers. 50 stochastic passes, epistemic var = variance of means.
    \item \textbf{``Key hyperparameters?''} --- $L\!=\!96$, patch$\!=\!16$, stride$\!=\!8$, $d\!=\!32$, depth$\!=\!2$, dropout$\!=\!0.3$, lr$\!=\!10^{-4}$, batch$\!=\!32$, mask ratio$\!=\!0.4$, MC samples$\!=\!50$.
    \item \textbf{``How long does training take?''} --- $\sim$3--5 min per experiment on one GPU. All 14 in $\sim$1 hour via \texttt{run\_all.sh}.
    \item \textbf{``What extra packages?''} --- Only \texttt{matplotlib}. Core: torch, numpy, pandas are in the base environment.
    \item \textbf{``Why $\lambda = 0.3$ for aux MSE?''} --- Tested \{0.1, 0.3, 0.5\}; 0.3 is the best MSE-vs-calibration trade-off.
    \item \textbf{``Show us a code segment?''} --- (Have \texttt{train.py} open at line 593 for the \texttt{ProbabilisticTSLANet} class, and line 688 for \texttt{GaussianNLLLoss}.)
\end{itemize}
\end{speakerbox}

\end{document}
